\section{Key Principles of Local-First Software}
\label{sec:key-principles-of-local-first-software}

Local-first software combines the independence of applications that store data locally with the synchronization and collaboration offered by cloud applications.
The key architectural difference is which copy of the data is treated as authoritative.
In a cloud application, the server usually holds the authoritative copy, while data stored on a client is a cache.
In a local-first application, the copy on the user's device is primary: the user can read and modify it without waiting for a server.
Servers can still store additional copies and help devices exchange changes, but the user's ability to work does not depend on them.
\parencite{local-first-software}

\textcite{local-first-software} describe seven principles for this approach: responsiveness, access across devices, offline operation, collaboration, longevity, privacy and security, and user ownership and control.
These principles concern both software architecture and the user's relationship with their data.
They serve as criteria for evaluating an application, and adding local storage or a synchronization library does not satisfy them automatically.

\subsection{Responsiveness through local operations}

The first principle is that users should not have to wait for a network request before continuing their work.
Reading and writing local data removes network latency from the immediate interaction, while synchronization can take place in the background.
For example, editing a locally available document should update and save the local copy even when the synchronization server is slow or unreachable.
\parencite{local-first-software}

This differs from an optimistic user interface, which displays an expected result before a server confirms the operation.
Such an interface can hide network latency, but the apparent success may need to be reversed if the request fails.
In a local-first application, a locally saved edit does not depend on a successful server request.
Local execution still does not guarantee good performance, since processing, storage access, and rendering take time.
The principle removes a remote dependency from ordinary interactions, which makes near-instant responses possible.
\parencite{local-first-software}

\subsection{Access across multiple devices}

Local storage should not limit a user's work to one computer.
The second principle requires synchronization across the devices on which the user works, such as a laptop, tablet, and phone.
Each device should be able to maintain a local copy of the relevant data and exchange changes with the others.
A server can act as an additional replica, keeping data available for synchronization when other devices are offline.
\parencite{local-first-software}

Synchronization introduces requirements beyond copying files.
If a user edits the same document on two disconnected devices, both copies may change before either device receives the other's updates.
The application must account for these concurrent changes when the devices reconnect.
An application that stores data locally but cannot synchronize it across devices does not meet this principle.

\subsection{Independence from network availability}

The third principle makes the network optional for working with data already available on the device.
Users should be able to both read and write while offline, with changes synchronized when a connection becomes available.
Offline operation is treated as a normal mode of use in which editing remains possible.
Devices can also synchronize without an Internet connection, for example over a local network or Bluetooth.
\parencite{local-first-software}

Network independence also requires access to the software needed to interpret and modify the data.
For a web application, the corresponding requirement is to make the application itself, as well as its data, available without fetching essential resources from a remote server.
Disconnected users still receive other people's changes only after they reconnect, but the missing changes do not stop them from continuing their own work.

\subsection{Collaboration across independent copies}

Local-first software aims to offer collaboration comparable to that of cloud applications while allowing participants to work independently.
When connected, users should be able to see one another's changes in real time.
When disconnected, they should be able to continue editing and combine their work later.
This requires a way to reconcile concurrent changes without making a central server a prerequisite for every edit.
\parencite{local-first-software}

Collaboration also extends beyond simultaneous editing of a single shared version.
Users may need to prepare changes privately, propose them for review, and selectively accept another person's contributions.
The original principles include both immediate collaboration and workflows similar to suggested edits or pull requests.
A collaborative system therefore has to help people review and understand changes, in addition to delivering those changes to other devices.
\parencite{local-first-software}

\subsection{Long-term access to data}

The fifth principle is that users should retain access to their data even after the service provider ceases to operate.
A user's work should remain accessible if the company that developed the application closes or stops maintaining its servers.
Storing both the data and the software needed to read and modify it locally reduces this dependency.
It also lets users keep a working version instead of being forced to accept every service-side upgrade.
\parencite{local-first-software}

Longevity is broader than surviving a short network outage.
An application might support several days of offline editing yet still depend on a provider to restore access or continue operating in the longer term.
Meeting this principle requires considering what remains usable after the provider's infrastructure disappears.
A local copy offers limited independence if the user cannot interpret or modify it without that infrastructure.

\subsection{Security and privacy by default}

The sixth principle aims to reduce the trust users must place in service providers.
Storing data on users' devices can reduce the need for a central database where the provider can read everyone's data.
Where servers are used for synchronization or backup, end-to-end encryption can allow them to store and relay encrypted data without being able to read its contents.
End-to-end encryption means that only the participants' devices hold the keys, so a server cannot decrypt the data it relays.
\parencite{local-first-software}

These properties require explicit design choices.
Local storage alone does not encrypt data, and using a synchronization server does not imply that the server is unable to inspect it.
Access control is also more difficult when collaborators possess independent copies.
Preventing future access cannot by itself remove information that another participant has already received.
The original essay leaves the meaning of revoking access in this setting as an open question.
\parencite{local-first-software}

\subsection{User ownership and control}

The final principle defines ownership in terms of practical agency rather than intellectual property law.
Users should be able to copy, modify, back up, and manage their data without depending on a provider's permission.
Possession of a local copy is part of this control, but the application must also avoid artificial restrictions on what users can do with it.
\parencite{local-first-software}

The authors do not make open-source licensing a requirement: commercial, closed-source software can satisfy the principles if it preserves this freedom.
Control also brings responsibility.
Users need ways to protect their data against device loss, accidental deletion, and ransomware, including maintaining backups.
\parencite{local-first-software}

\subsection{Architectural implications and limits}

Taken together, the principles require local persistence, synchronization between independently editable copies, and continued access beyond a particular service.
Servers remain useful for discovery, backup, and exchanging changes between devices that are not online at the same time.
What distinguishes the architecture is that losing contact with a server does not prevent users from working with their local data.
\parencite{local-first-software}

Conflict-free replicated data types (CRDTs) provide one way to reconcile independently edited copies.
They define data structures and merge rules that allow replicas receiving the same updates to converge to the same state.
For example, the JSON CRDT described by \textcite{json-crdt} supports nested maps and lists and performs merging on the clients without depending on network message ordering.
Such mechanisms support offline editing and collaboration, but do not by themselves provide transport, access control, or long-term usability.

Convergence also does not guarantee that a merged document expresses its authors' intentions.
The Upwelling research prototype shows this with concurrent replacements of a word: a text merge can preserve both replacements as a combined string that still requires human correction.
Review workflows and understandable change history therefore remain important even when replicas converge automatically.
\parencite{upwelling}

An application is best assessed against all seven principles at once.
Offline support alone says little about privacy or long-term control, and local file ownership alone does not provide seamless collaboration.
Local-first software aims to provide networked collaboration while keeping users independent of the service that supports it.
\parencite{local-first-software}
